
\section{Twenty-Third Sunday after Trinity: Then surely the children are free}

\begin{quote}

Matthew 17:24–27. And when they had come to Capernaum, those who collected the Temple tax came to Peter and said, Does your Master not pay the Temple tax? He said, Yes. And when he had come into the house, Jesus anticipated him and said, What do you think, Simon? From whom do the kings of the earth take toll or tax, from their own children or from strangers? Peter says to him, From strangers. Jesus said to him, Then surely the children are free. But so that we may not give offense to them, go to the sea, cast a hook, and take the first fish that comes up; and when you open its mouth, you will find a stater; take that and give it to them for me and you.

\end{quote}

\bigskip

The solemn lesson which is contained in our text today Paul has expressed thus in the Epistle to the Galatians: "You were called to freedom, brothers! Only do not use freedom as an occasion for the flesh, but serve one another in love."

Jesus had just come home to Capernaum from the journey during which Simon Peter had confessed concerning him: "You are the Christ, the Son of the living God;" and immediately after this confession three of the disciples had been witnesses of his Transfiguration and had heard the voice from the cloud: "This is my beloved Son, in whom I am well pleased; hear him." It had therefore been witnessed on earth and witnessed from heaven that Jesus was the Son, the only-begotten of the Father, exalted above all and all, so that he was in truth Lord over all things, and no one was lord or ruler over him. If there were ever to be any talk of someone being free, then surely it must be Jesus, the King’s Son, whom all ought to honor, the Teacher of truth, whom all ought to hear.

But those who collected the Temple tax knew nothing of Jesus’ infinite majesty and did not understand that he is Lord of all. In narrow official pride and Jewish zeal for the old ordinances they pressed on with their tax collecting. It never entered their mind that there could be any exception to the rule, or that there could be any talk of any change in the "divine" ordinance, as they would call it. Perhaps these tax collectors were doubly eager to get the tax from Jesus because they had something like a feeling that he held quite another view of the old ordinances than they did, and that they might perhaps get occasion to accuse him of breaking the Law if he should refuse to pay the Temple tax. They met Peter and asked him, with some reproach, it seems: "Does your Master not pay the Temple tax?" They thought it ought to have been paid already long ago, and that it suited such a man poorly to be in arrears with so holy an obligation.

And Peter hastened to answer: Yes. All that he had seen and heard from Jesus and about Jesus in these last weeks had by no means driven self-will and the fear of man out of Peter. He still thought that he understood how matters ought to be arranged, and how Jesus ought to avoid "suffering many things from the elders and chief priests and scribes." He did not want any misunderstanding with the authorities and to expose Jesus to any accusation of rebellion or lawbreaking. He too had understood that Jesus thought otherwise than others about the old ordinances, and he wished to avoid all unpleasantness by quickly settling the matter on his own. Should the tax collectors meet Jesus himself, then in Peter’s judgment there might easily be another clash, as there had just lately been precisely over the ordinances of the elders (Matt. 15:1 and 2).

But Peter had acted too hastily when he took the decision in this matter into his own hand. And yet we can scarcely reproach him for it when we read of the wise and loving way in which Jesus set him right. What should we not have lost from the Gospel if the disciples’ missteps had not been the occasion for Jesus to reveal his heart! Peter was given no time to explain the matter in his own way; Jesus anticipated him, for he would surely not have Peter become more deeply entangled in his thoughtlessness. And Jesus makes the matter clear with a single word: "From whom do the kings of the earth take toll or tax, from their own children or from strangers?" Peter answers: "From strangers." And Jesus’ conclusion is self-evident: Then surely the children are free.

Yes, this is the main thing: the children are free. If we truly and indeed are God’s children, then it is altogether certain that freedom is ours. The learned have disputed much as to whether the children here in this place are only Jesus, or whether it is Jesus and his disciples. But if they had understood what is written, that the one whom the Son sets free is truly free, then they would not have disputed about this. For however great the difference is between the Father’s only-begotten Son and God’s children, yet this is true, that all that the Son has received from the Father, that he has given to his believers, according to what he himself says: "And I have given them the glory that you have given me." God’s children are God’s heirs and Jesus Christ’s fellow heirs, so that the same freedom that he has, his brothers and sisters also have.

The children are free. And there are surely few words that contain so much and that are so difficult to grasp. From the compulsion of the Law and from rules and ordinances the children are free, because the love of God is poured out in their hearts through the Holy Spirit, who is given to them.

But if we feel with inward shame today that we are not free, that we are dragged along with so much sin, we are dragged along with so much guilt, we are plagued with so many rules and forms and bonds, then it is still time to win the children’s freedom. Come to Jesus’ cross and see what sin and the Law and wrath and death have done with the Son. Do you truly believe that he suffered it all for you? Do you truly believe it? Very well then, where is your sin and guilt, where is God’s wrath and judgment over you? Are you not free when the Lamb of God has taken it all upon himself? How long will you torment yourself with a yoke that Jesus will take altogether away, if only you will entrust yourself wholly to him? The children are indeed free. There is no guilt of sin and no sentence of death over those who have placed their cause into Jesus’ hand.

Come to Jesus and see how he goes so freely and boldly in the service of love without heeding commandments and prohibitions. Only look at him, how little he thinks of becoming unclean when he touches a leper, or when he takes the dead girl by the hand and raises her up! Look at him, how free he is when he heals on the Sabbath and commands the healed man: "Take up your bed and walk!" Learn from him that love is free, and mercy glories over judgment, and do not let yourself be bound by the many thousands of human considerations and rules.

The children are free. And if you have not yet learned this, then come to him who had nowhere to lay his head, and learn from him how free the one is who can say with Paul: "I have learned to be content with what I have; I know how to be brought low, and I know also how to abound; in everything and in all things I am instructed, both to be full and to hunger, both to abound and to suffer need; I can do all things in Christ, who strengthens me." Yes, truly in all things the children are free. God is their Father, and he loves them; therefore all his ways are blessed ways of freedom for them. Jesus is their brother, and in him they believe; therefore they are altogether certain of the forgiveness of their sins, and Christ’s atonement is their charter of freedom from the Law’s accusation. The Holy Spirit is poured out in their hearts; therefore they are moved willingly in the ways of love. Truly the children are free.

Therefore Jesus did not need to pay the Temple tax either. Yet he adds: "But so that we may not give offense to them, go to the sea, cast a hook, and take the first fish that comes up! And when you open its mouth, you will find a stater; take that and give it to them for me and you."

See again the wisdom and the love! He who has all God’s power and all God’s freedom does not use it for his own exaltation and his opponents’ humiliation. For him it matters not to give offense to them. That they tried to give offense to him, that was their fault and their misfortune; he will not give offense to them by doing what he had a right to do, but what they were unfit to understand. Jesus did not yield to the Pharisees when it was a matter of healing on the Sabbath and sparing the sick person a day of suffering; but he did not use his right and freedom when it was he himself who had to suffer by the sacrifice. He let the tax collectors have the tax they demanded, and did not reveal his power and his dominion and his freedom to them.

But Peter was to try and see that the children are free even then, when love compels them to hide their freedom. Therefore the Savior sent him out to fish and gave him this time also a wondrous catch; for in the fish’s mouth Peter found a coin large enough to pay the Temple tax for both Jesus and Peter. This was to teach Peter that love should never be at a loss for means when it takes burdens and difficulties upon itself "so as not to give offense to them."

Thus the children are again free, for God is himself the one who gives them all that they use in the service of love. The fish and the stater in its mouth, all is God’s gift, and we do not give ours, but his. It is as David says when he and his people brought rich gifts for the building of the Temple: "Who am I, and what is my people, that we should have strength to give a freewill offering such as this? From you is it all, and from your hand we have given to you" (1 Chron. 29:14).

Thus there is no estate so glorious as that of God’s children. In freedom and love they walk through the world, and if they cannot spread light and joy upon their fellow men’s path, this at least is granted them: not to give offense.

May the Lord help us rightly to understand it: "You were called to freedom, brothers! Only do not use freedom as an occasion for the flesh, but serve one another in love" (Gal. 5:13).
